\section{Der gute Hirte}

\begin{block}[Einführung]
    Heute will ich euch Unterricht über Schafe erteilen und zwar über das Walliser Schwarznasenschaf. Das Walliser Schwarznasenschaf ist eine alte Schafrasse, die in den Alpen beheimatet ist. Sie ist bekannt für ihre Widerstandsfähigkeit und Anpassungsfähigkeit an die rauen Bedingungen der Bergwelt.

    Diese Schafe sind soziale Tiere und leben in Herden zusammen. Es gibt aber auch Einzelgänger, aber eher selten. Die Schafe fressen zusammen, haben dabei die Köpfe zusammen um so besser mitzubekommen, was sein Nachbar frisst und ob es sich lohnt ihm das wegzufressen.

    Über die warme Tageszeit von 10 - 17Uhr stecken sie ihre Nasen irgendwo unter einen Stein in den Schatten, kauen und warten bis es wieder kühler wird.

    Wenn der Hirte ruft, gucken alle zu ihm hin und wenn das erste Schaf losläuft, laufen alle anderen hinterher. Die laufen nicht weil der Hirte so nett ruft, sondern weil sie denken, dass sie etwas Gutes bekommen. Legt man ihnen regelmässig etwas Leckeres aus, wie Salz oder Gläck, dann kommen sie von alleine um zu schauen, ob da wieder was liegt. Interessant ist, sie vergessen über Jahre nicht, wo diese speziellen Stellen sind. Ein Jahr später, wenn sie wieder in der Gegend sind, gehen sie genau zu diesen Stellen und schauen ob da noch was liegt. Finden sie nichts gehen sie wieder.

    Kommt euch dieses Verhalten nicht ein bisschen bekannt vor? Sind wir Menschen nicht auch ein bisschen so? Mal gucken wir, wie der Garten des Nachbarn aussieht, ist ja gewöhnlich immer grüner als der Eigene. Welches Auto der so hat, wer im Urlaub brauner wurde. Der Mensch ist nicht gerne alleine. Wenn auf dem Bahnhofplatz einer steht und in die Luft guckt, dann gucken alle anderen auch in die Luft. Wenn ein Geschäft Aktionen hat, dann rennen alle dorthin und gehen später nochmals hin, um zu schauen, ob es noch was gibt.
    
    Der Mensch will Gesellschaft, obwohl er eigentlich egoistisch ist. Er will alle Freiheiten und um diese zu erhalten, hat er kein Problem, diese Freiheit seinem Mitmenschen wegzunehmen. Der Mensch selbst ist gewissenlos. Ohne Gesetze und Regeln funktioniert ein Zusammenleben nicht.

    Genauso wie das Walliser Schwarznasenschaf ohne Hirte oder Züchter nicht überleben kann, funktioniert ein Leben bei uns Menschen ohne Hirte auch nicht. Wir brauchen jemanden, der uns führt und das seit dem Sündenfall im Paradies.

    \begin{block}[Gott der Hirte]
        Vielleicht habt ihr es schon erraten, heute geht es um Namen Gottes als \betonung[b, p]{Hirte oder guter Hirte.} Gott hat sich ein Volk auserwählt und diesem Volk war er ein Hirte, ein Führer, ein Leiter. Dieses Volk sollte unter der Führung Gottes, den umliegenden Völkern ein gutes Beispiel sein, so dass sich die Völker auch zu diesem Hirten zuwenden.
        \biblezit{Jes}{40:11}

        Schon sehr früh zeigte sich, dass das Volk nicht auf ihren Hirten hören wollte. Der Hirte hatte ihnen versprochen, wenn sie auf ihn hören, dann werde es ihnen gut gehen. Er würde dafür sorgen, dass sie immer genug zu Essen und zu Drinken haben und in Frieden leben können

        Der \biblezit{Ps}{23:1-3} zeigt sehr schön wie Gott sein Volk als Hirte führt. Wir finden im alten Testament noch andere Stellen, in denen Gott seinem Volk zeigt, wie er sie unterstützen und führen wollte.
        \biblezit{Ps}{80:2}, \biblezit{Ps}{100:3}

        Als sein Sohn Jesus Christus auf diese Welt kam, um sein Erlösungwerk zu vollenden, hat auch Jesus immer wieder dieses Bild vom Hirten benutzt.

        Aber das Volk wollte nicht. Sie wollten keinen Hirten, der ihnen den Weg zeigt, sie beschützt und das Tor für das ewige Leben öffnet. \biblezit{Sach}{10:2} Jesus, der gute Hirte, wurde geschlachtet wie ein Lamm. Als makelloses Lamm ist Jesus für unsere Sünden ans Kreuz gegangen und hat für uns gelitten und ist gestorben. Er, der gute Hirte, ist für seine Herde gestorben. 

        Am Ende der Zeit im tausendjährigen Reich wird Jesus wiederkommen. Dieses Mal aber in Macht und Herrlichkeit und wird nicht mehr nur ein Hirte über sein Volk sein, sondern der König und Herrscher.
        \biblezit{Sach}{14:9}

        Im ersten Petrusbrief \zB{} \biblezit{IPetr}{5:4} beschreibt Petrus Jesus sogar als den obersten Hirten. 
        Petrus hat das Evangelium vom Menschen Jesus, als er auf der Erde war, erhalten, wo Jesus von sich selbst als vom guten Hirten spricht. \biblezit{Joh}{10:11-16}

        Paulus hat sein Evangelium vom erhöhten Jesus erhalten, nach seiner Auferstehung.
        Paulus redet nach der Auferstehung von Jesus Christus nicht mehr darüber, dass Jesus unser Hirte ist. Er sagt aber, dass der, der das Erlösungswerk unseres Herrn Jesus annimmt und von Herzen glaubt, ein Erbe ist, der die Sohnschaft empfangen hat. Der ist nun ein Teil von der Gemeinde, die als Körper mit dem Haupt Jesus Christus verbunden wird. \biblezit{IKor}{12:12-27} \biblezit{Gal}{4:4-7}


        Klar dürfen wir immer noch das Bild vom Hirten und den Schafen auch für uns in Anspruch nehmen. Es ist ein schöner Vergleich, wie wir allein, ohne diesen Hirten, in dieser Welt untergehen. In einer Welt, in der Satan wie ein brüllender Löwe umhergeht, \biblezit{IPetr}{5:8} gibt uns der Psalm 23 immer wieder Trost, wenn wir wissen, dass wir uns im Herrn geborgen fühlen können, dass er uns nur soviel aufbürdet, wie wir auch tragen können. Ich weiss für jemanden, dem es schlecht geht, ist es schwierig in dieser Zeit darin einen gütigen Gott zu erkennen. Aber was haben wir sonst? Die Krankheit oder die Probleme gehen auch nicht weg, wenn wir mit Gott hadern und ihn verleugnen. Im Gegenteil, es wird nur noch schlimmer. Ist es nicht ein Trost, zu wissen, dass es einmal besser wird? Und das nicht nur für eine kurze Zeit? Wie hier auf Erden wo die Ferien auf der Alpe gerade mal 4 Wochen dauerten? Sondern für die Ewigkeit!

        Bis dahin finden wir Trost im Glauben, im Gebet, in der Gemeinschaft der Gemeinde. Nutzt die Gemeinde für Hilfe und Unterstützung, dafür hat Gott sie so eingesetzt, wie sie in der Bibel beschrieben ist, dass jeder mit seinen spezifischen Gaben den jeweiligen Brüdern und Schwestern hilft. Das hilft nicht nur uns, sondern dadurch geben wir auch gegen aussen ein gutes christliches Bild ab, das wiederum unseren Herrn und Hirten ehrt und lobt.

    \end{block}


\end{block}